\subsection{一致空间的积}

定义 4 若 \((\mathbf{X}_t)_{i\in I}\) 是一族一致空间，则该族的积一致空间是指积集合

\[
\mathbf {X} = \prod_ {i \in I} \mathbf {X} _ {i}
\]

并赋以使诸投影 \(pr_{t}: X \rightarrow X_{t}\) 都一致连续的最粗一致结构。这一一致结构称为诸 \(X_{t}\) 的一致结构的积，而一致空间 \(X_{t}\) 称为 X 的因子。

积一致结构在 X 上诱导的拓扑与诸 \(X_{i}\) 的拓扑的积相同（第 3 目命题 4 的推论）。

命题 7 设 \(f = (f_{t})\) 是一致空间 \(Y\) 到积一致空间 \(X = \prod_{i \in I} X_{i}\) 的映射。则 \(f\) 一致连续当且仅当每个 \(f_{t}\) 都一致连续。

由于 \(f_{i} = \mathrm{pr}_{i} \circ f\)，这是第 3 目命题 4 的一个特殊情形。

推论 设 \((\mathbf{X}_t)_{i \in I}, (\mathbf{Y}_t)_{i \in I}\) 是以同一集 I 为指标的两族一致空间。对每个 \(i \in I\)，设 \(f_t\) 是 \(\mathbf{X}_t\) 到 \(\mathbf{Y}_t\) 的映射。若每个 \(f_t\) 都一致连续，则积映射

\[
f: (x _ {i}) \rightarrow (f _ {i} (x _ {i})).
\]

反之，若诸 \(X_{i}\) 非空且 f 一致连续，则每个 \(f_{i}\) 都一致连续。

\(f\) 可写成 \(x \to (f_{i}(\mathrm{pr}_{i}x))\)，因此推论的第一部分由命题 7 得出。第二部分的证明：考虑一点 \(a = (a_{i})\)（属于 \(\prod_{i \in I} X_{i}\)），并重复第一章 § 4 第 1 目命题 1 推论 1 的论证，把其中“在 \(a\)（相应地，\(a_{x}\)）连续”一语换成“一致连续”。

初始一致结构的可递性的一般判据（第 3 目命题 5）表明，如同拓扑空间的积的情形一样（第一章 § 4 第 1 目），一致空间的积是可结合的，并且下述结论成立：

命题 8 设 X 是一个集合，\((\mathbf{Y}_t)_{t \in I}\) 是一族一致空间，且对每个 \(t \in I\)，\(f_t\) 是 X 到 \(\mathbf{Y}_t\) 的映射。设 \(f\) 是映射 \(x \to (f_t(x))\)（从 X 到 \(\mathbf{Y} = \prod_{t \in I} \mathbf{Y}_t\)），并设 \(\mathfrak{U}\) 是 X 上使诸 \(f_t\) 一致连续的最粗一致结构。则 \(\mathfrak{U}\) 是在 \(f\) 之下，Y 的积一致结构在 \(f(\mathbf{X})\) 上诱导的一致结构的逆像。

推论 对每个 \(i \in I\)，设 \(A_{i}\) 是 \(Y_{i}\) 的子空间。则在 \(A = \prod_{i \in I} A_{i}\) 上由 \(\prod_{i \in I} Y_{i}\) 的积一致结构所诱导的一致结构，与诸子空间 \(A_{i}\) 的一致结构的积相同。

此外，立即可以看出，若 \(X_{1}, X_{2}\) 是两个一致空间而 \(a_{1}\) 是 \(X_{1}\) 的任意一点，则映射 \(x_{2} \to (a_{1}, x_{2})\) 是 \(X_{2}\) 到子空间 \(\{a_{1}\} \times X_{2}\)（\(X_{1} \times X_{2}\) 的）上的同构；因此：

命题 9 设 \(f\) 是积一致空间 \(\mathbf{X}_1 \times \mathbf{X}_2\) 到一致空间 \(\mathbf{Y}\) 的一致连续映射；则每个偏映射

\[
x _ {2} \rightarrow f (x _ {1}, x _ {2})
\]

（作为 \(X_{2}\) 到 Y 的映射）都是一致连续的。

换言之，两个变元的一致连续函数对每个变元分别是一致连续的。

* 第一章 § 4 第 2 目注 2 中给出的例子表明，本命题的逆命题不成立。 *

